\section{Akzeptanzfaktoren im organisationalen Kontext}
\label{sec:05-05_acceptance-factors}

Die Akzeptanz des neuen Wikis im organisationalen Kontext wurde sowohl über geschlossene Skalenfragen als auch über qualitative Interviewaussagen erhoben. Im Fokus standen dabei Nutzungsmotivation, tatsächliche Nutzungsintention sowie fördernde und hemmende organisationale Rahmenbedingungen. Die Ergebnisse zeigen, dass die Akzeptanz des neuen Wikis weniger von individuellen Faktoren als vielmehr von strukturellen und kommunikativen Bedingungen innerhalb der Organisation geprägt ist.

Die geschlossene Frage zur zukünftigen Nutzungsintention (GF-1) wurde überwiegend positiv beantwortet. Der arithmetische Mittelwert lag bei 4{,}1 auf einer fünfstufigen Skala, der Median bei 4, bei einer Standardabweichung von 0{,}7. Sechs Teilnehmende gaben an, das neue Wiki künftig regelmäßig nutzen zu wollen, drei Personen wählten eine eher zustimmende Antwortoption und eine Person äußerte eine neutrale Haltung; ablehnende Bewertungen traten nicht auf (vgl.~Anhang). Damit zeigt sich eine insgesamt hohe Nutzungsbereitschaft, bei gleichzeitig etwas größerer Streuung als bei den Bewertungen zur Informationsarchitektur oder Rollenpassung.

In den qualitativen Begründungen zu dieser Frage (HF-4) wurden mehrere Akzeptanztreiber identifiziert. Sechs von zehn Teilnehmenden nannten aktuelle, verlässliche Inhalte als zentrale Voraussetzung für eine dauerhafte Nutzung. Fünf Personen betonten die Bedeutung des Wikis als zentrale Wissensquelle, um parallele Informationsablagen zu vermeiden. Ebenfalls fünf Teilnehmende hoben hervor, dass eine klare Struktur und gute Auffindbarkeit ihre Bereitschaft zur Nutzung erhöhen. Vier Personen verwiesen auf den konkreten Nutzen im Arbeitsalltag, etwa zur schnellen Klärung von Zuständigkeiten oder zur Vorbereitung von Terminen (vgl.~Anhang).

Gleichzeitig wurden in den Interviews mehrere Akzeptanzbarrieren benannt. Sechs Teilnehmende äußerten die Sorge, dass veraltete oder unvollständige Inhalte die Nutzung langfristig hemmen könnten. Fünf Personen verwiesen auf die eingeschränkte Leistungsfähigkeit der Suchfunktion als potenzielles Hindernis. Vier Teilnehmende nannten das Nebeneinander mehrerer Wissensquellen als problematisch, da dies zu Unsicherheit darüber führen könne, welche Informationen als verbindlich gelten. Diese Barrieren beziehen sich überwiegend auf organisationale Prozesse und Zuständigkeiten und nicht auf die grundsätzliche Gestaltung des Wikis (vgl.~Anhang).

Ein weiterer zentraler Akzeptanzfaktor betrifft die begleitende Kommunikation und Einführung des Wikis. Acht von zehn Teilnehmenden gaben an, dass eine aktive Vorstellung des Wikis in bestehenden Meetings oder Gremien die Nutzung fördern würde. Sechs Personen wünschten sich dialogische Formate wie Fragerunden oder kurze Austauschformate, um Unsicherheiten zu klären und Nutzungsszenarien zu diskutieren. Ebenfalls sechs Teilnehmende nannten regelmäßige Hinweise oder Updates, etwa in Form von Newslettern, als unterstützend, betonten jedoch, dass diese allein nicht ausreichen würden. Fünf Personen hoben die Rolle von Führungskräften als Vorbilder und Multiplikatoren hervor (vgl.~Anhang).

Die offenen Feedbackfragen (OF-1 bis OF-4) bestätigen dieses Bild. Mehrere Teilnehmende äußerten explizit, dass die langfristige Akzeptanz weniger von weiteren funktionalen Erweiterungen als von klar geregelten Verantwortlichkeiten für Pflege und Aktualisierung abhänge. Verbesserungsvorschläge bezogen sich häufig auf organisatorische Maßnahmen, etwa klare Zuständigkeiten für Inhalte, transparente Kommunikationsregeln oder verbindliche Nutzungserwartungen. Einzelne Rückmeldungen betrafen zusätzliche Schulungs- oder Onboarding-Angebote, insbesondere für neue oder externe Mitarbeitende (vgl.~Anhang).

Zusammenfassend zeigen die Ergebnisse, dass das neue Wiki auf individueller Ebene auf hohe Akzeptanz und Nutzungsbereitschaft stößt. Die entscheidenden Einflussfaktoren für eine nachhaltige Nutzung liegen jedoch primär im organisationalen Kontext. Insbesondere die Sicherstellung aktueller Inhalte, die Reduktion paralleler Wissensquellen sowie eine aktive, dialogische Einführung und Kommunikation werden von der Mehrheit der Teilnehmenden als zentrale Voraussetzungen für Akzeptanz benannt. Diese Befunde schließen die Ergebnisdarstellung ab und bilden die Grundlage für die anschließende Diskussion und Einordnung im folgenden Kapitel.